\chapter{میانه‌روی: هنر یا بی‌اصالتی؟}

\begin{tcolorbox}[
    enhanced,
    colback=VeryLightGray, % FIXED: use defined color name
    colframe=CenterGreen, % FIXED: use defined color name
    boxrule=2pt,
    arc=8pt,
    fonttitle=\bfseries\large,
    title={\rl{نقل‌قول آغازین}},
    attach boxed title to top center={yshift=-3mm},
    boxed title style={colback=CenterGreen, colframe=CenterGreen} % FIXED: use defined color name
]
\begin{center}
\Large\textit{«فضیلت در میانه است.»}

\vspace{3mm}
\normalsize
--- ارسطو، \textit{اخلاق نیکوماخوس}، قرن چهارم پیش از میلاد
\end{center}
\end{tcolorbox}

\vspace{5mm}

\section*{درآمد}


در سراسر این کتاب از «چپ» و «راست» سخن گفته‌ایم. اما عنوان کتاب پرسشی دیگر طرح می‌کند: «میانه کجاست؟» آیا میانه نقطه‌ای ثابت میان دو قطب است؟ آیا روشی برای تفکر است؟ یا صرفاً بی‌تصمیمی و فرصت‌طلبی؟

میانه‌روی در تاریخ اندیشه جایگاه والایی داشته—از «حد وسط» ارسطو تا «اعتدال» در سنت‌های دینی. اما در سیاست معاصر، میانه‌روی هم ستایش می‌شود (به عنوان عقلانیت و واقع‌گرایی) و هم نکوهش (به عنوان بی‌اصالتی و حفظ وضع موجود). این فصل به کاوش در معنا، تاریخ، و آیندهٔ میانه‌روی می‌پردازد.


%═══════════════════════════════════════════════════════════════
\section{میانه چیست؟}
%═══════════════════════════════════════════════════════════════

\subsection{سه معنای میانه}

\begin{center}
\begin{tikzpicture}[scale=1, every node/.style={font=\small}]

% Three boxes for three meanings
\node[draw=centergreen, fill=centergreen!10, rounded corners, minimum width=4cm, minimum height=3cm, align=center] at (0,0) {
    \textbf{\large\rl{۱. میانه به مثابه نقطه}}\\[3mm]
    \rl{موضعی بین چپ و راست}\\
    \rl{مثلاً: بین سوسیالیسم}\\
    \rl{و سرمایه‌داری ناب}
};

\node[draw=rightblue, fill=rightblue!10, rounded corners, minimum width=4cm, minimum height=3cm, align=center] at (5.5,0) {
    \textbf{\large\rl{۲. میانه به مثابه روش}}\\[3mm]
    \rl{پراگماتیسم و انعطاف}\\
    \rl{رد ایدئولوژی سخت}\\
    \rl{«هر چه کار می‌کند»}
};

\node[draw=goldaccent, fill=goldaccent!10, rounded corners, minimum width=4cm, minimum height=3cm, align=center] at (11,0) {
    \textbf{\large\rl{۳. میانه به مثابه فضا}}\\[3mm]
    \rl{جایی که اکثریت هستند}\\
    \rl{متغیر با زمان و مکان}\\
    \rl{«افکار عمومی میانه»}
};

% Title
\node[above, font=\bfseries\large] at (5.5,2.5) {\rl{سه معنای «میانه» در سیاست}};

\end{tikzpicture}
\end{center}


\textbf{۱. میانه به مثابه نقطه:}
در این معنا، میانه موضعی محتوایی است: نه سوسیالیسم کامل، نه سرمایه‌داری ناب؛ نه آزادی مطلق، نه کنترل مطلق. این نگاه فرض می‌کند که طیف سیاسی خطی است و میانه نقطه‌ای میان دو سر است.

\textbf{۲. میانه به مثابه روش:}
در این معنا، میانه‌روی نه موضع بلکه رویکرد است: پراگماتیسم، تجربه‌گرایی، انعطاف‌پذیری. میانه‌رو کسی است که به ایدئولوژی سخت پایبند نیست و بر اساس شواهد تصمیم می‌گیرد.

\textbf{۳. میانه به مثابه فضای اجتماعی:}
در این معنا، میانه جایی است که اکثریت مردم قرار دارند. این میانه ثابت نیست؛ با تغییر افکار عمومی جابجا می‌شود. آنچه امروز میانه است، ممکن است فردا چپ یا راست باشد.

هر سه معنا مشکلاتی دارند. معنای اول فرض می‌کند که «حق» همیشه در وسط است—که لزوماً درست نیست. معنای دوم می‌تواند به بی‌اصولی بینجامد. معنای سوم میانه را نسبی و متغیر می‌کند.


%═══════════════════════════════════════════════════════════════
\section{فلسفهٔ میانه‌روی}
%═══════════════════════════════════════════════════════════════

\subsection{ارسطو و حد وسط}

\begin{tcolorbox}[
    enhanced,
    colback=goldaccent!10,
    colframe=goldaccent,
    boxrule=1.5pt,
    arc=5pt,
    fonttitle=\bfseries,
    title={\rl{🟡 زمینهٔ فلسفی: نظریهٔ حد وسط ارسطو}}
]
\textbf{ارسطو} (۳۸۴-۳۲۲ پ.م.) در \textit{اخلاق نیکوماخوس} نظریهٔ «حد وسط» (Mesotes) را مطرح کرد. بر اساس این نظریه، فضیلت اخلاقی در میانهٔ دو افراط قرار دارد:

\begin{center}
\begin{tabular}{ccc}
\textbf{کمبود} & \textbf{فضیلت (حد وسط)} & \textbf{افراط} \\
\hline
ترسو & شجاع & بی‌پروا \\
خسیس & سخاوتمند & ولخرج \\
منفعل & خویشتن‌دار & عصبانی \\
\end{tabular}
\end{center}

اما ارسطو تأکید می‌کرد که حد وسط «میانگین ریاضی» نیست. حد وسط برای هر فرد و در هر موقعیتی متفاوت است و یافتن آن نیازمند \textbf{حکمت عملی} (Phronesis) است.
\end{tcolorbox}


آیا حد وسط ارسطویی به سیاست قابل تعمیم است؟ ارسطو خود در \textit{سیاست} از «حکومت میانه» (Polity) دفاع می‌کرد—حکومتی که نه الیگارشی باشد نه دموکراسی محض، بلکه ترکیبی از هر دو.

ارسطو معتقد بود که طبقهٔ متوسط ستون ثبات سیاسی است. جامعه‌ای که در آن طبقهٔ متوسط قوی باشد، کمتر دچار انقلاب و بی‌ثباتی می‌شود.

این ایده‌ها در تاریخ اندیشهٔ سیاسی تأثیرگذار بوده‌اند. «حکومت مختلط» رومی، «موازنهٔ قوا» مونتسکیو، و حتی برخی نظریات دموکراسی معاصر ریشه در این سنت دارند.


\subsection{اعتدال در سنت‌های دینی}


مفهوم اعتدال در سنت‌های دینی مختلف حضور دارد:

\textbf{اسلام:}
قرآن مسلمانان را «امت وسط» (أُمَّةً وَسَطًا) می‌نامد. اعتدال (وسطیت) در فقه و اخلاق اسلامی اصلی مهم است. افراط و تفریط هر دو مذموم‌اند.

\textbf{بودیسم:}
«راه میانه» (Madhyamaka) یکی از آموزه‌های اصلی بودایی است. بودا راهی میان زهد افراطی و لذت‌جویی افراطی آموخت.

\textbf{کنفوسیوسیسم:}
\textit{ژونگ یونگ} (میانه‌روی) یکی از چهار کتاب کنفوسیوسی است. اعتدال در همه چیز—احساسات، رفتار، سیاست—آموخته می‌شود.

\textbf{مسیحیت:}
هرچند مسیحیت تأکید کمتری بر اعتدال دارد، سنت فلسفی مسیحی (توماس آکویناس) ارسطو را پذیرفت.

این سنت‌ها نشان می‌دهند که میانه‌روی نه اختراع مدرن، بلکه آرمانی کهن است.


\subsection{پراگماتیسم}

\begin{tcolorbox}[
    enhanced,
    colback=rightblue!8,
    colframe=rightblue,
    boxrule=1.5pt,
    arc=5pt,
    fonttitle=\bfseries,
    title={\rl{🔵 مفهوم کلیدی: پراگماتیسم}}
]
\textbf{پراگماتیسم} مکتب فلسفی آمریکایی است که حقیقت را بر اساس نتایج عملی می‌سنجد. بنیان‌گذاران: \textbf{ویلیام جیمز} (۱۸۴۲-۱۹۱۰) و \textbf{جان دیویی} (۱۸۵۹-۱۹۵۲).

\textbf{اصول پراگماتیسم:}
\begin{itemize}
    \item حقیقت آن چیزی است که «کار می‌کند»
    \item ایده‌ها ابزارند، نه تصویر واقعیت
    \item تجربه مهم‌تر از نظریهٔ پیشین است
    \item انعطاف‌پذیری فضیلت است
\end{itemize}

در سیاست، پراگماتیسم یعنی رد ایدئولوژی‌های سخت و تأکید بر آنچه در عمل نتیجه می‌دهد.
\end{tcolorbox}


پراگماتیسم سیاسی می‌گوید: مهم نیست سیاستی «چپ» است یا «راست»؛ مهم این است که کار می‌کند یا نه. این رویکرد با شعار تونی بلر («آنچه مهم است آن است که کار می‌کند») و شعار دنگ شیائوپینگ («مهم نیست گربه سیاه است یا سفید؛ مهم این است که موش بگیرد») بیان شده.

منتقدان می‌گویند پراگماتیسم می‌تواند پوششی برای بی‌اصولی باشد. «کار می‌کند» یعنی چه؟ برای چه کسی کار می‌کند؟ بر اساس چه معیاری؟

دفاع از پراگماتیسم: جهان پیچیده‌تر از آن است که یک ایدئولوژی بتواند همه چیز را توضیح دهد. تجربه و آزمون و خطا بهتر از جزم‌اندیشی است.


%═══════════════════════════════════════════════════════════════
\section{سانتریسم رادیکال؟}
%═══════════════════════════════════════════════════════════════


آیا می‌توان میانه‌رو بود و در عین حال رادیکال؟ برخی سیاستمداران و متفکران تلاش کرده‌اند «سانتریسم رادیکال» را تعریف کنند:

\textbf{ایده:} میانه‌روی لزوماً به معنای حفظ وضع موجود نیست. می‌توان از موضع میانه، اصلاحات بنیادین خواست—اصلاحاتی که نه ایدئولوژیک چپ باشند نه راست.

\textbf{نمونه‌ها:}
\begin{itemize}
    \item اصلاحات نظام انتخاباتی
    \item مبارزه با فساد
    \item شفافیت و پاسخگویی
    \item نوسازی نهادها
\end{itemize}

اما منتقدان می‌پرسند: آیا این واقعاً «میانه» است، یا فقط نام‌گذاری جدید برای مواضعی که می‌توانند چپ یا راست باشند؟


\subsection{ماکرونیسم: موردکاوی}

\begin{tcolorbox}[
    enhanced,
    colback=gray!5,
    colframe=darkgray,
    boxrule=2pt,
    arc=10pt,
    shadow={2mm}{-2mm}{0mm}{black!30},
    fonttitle=\bfseries\large,
    title={\rl{📋 موردکاوی: امانوئل ماکرون و «نه چپ، نه راست»}}
]

\textbf{امانوئل ماکرون} (متولد ۱۹۷۷) در ۲۰۱۷ با شعار «نه چپ، نه راست» به ریاست جمهوری فرانسه رسید.

\textbf{زمینه:}
\begin{itemize}
    \item بحران احزاب سنتی (سوسیالیست و جمهوری‌خواه)
    \item خستگی از تقسیم‌بندی قدیمی
    \item تهدید راست افراطی (لوپن)
\end{itemize}

\textbf{سیاست‌های ماکرون:}
\begin{itemize}
    \item اصلاحات بازار کار (راست)
    \item سرمایه‌گذاری در محیط زیست (چپ)
    \item لیبرالیسم اقتصادی (راست)
    \item حقوق اقلیت‌ها (چپ)
    \item اروپاگرایی قوی
\end{itemize}



\textbf{نقدها:}
\begin{itemize}
    \item «رئیس‌جمهور ثروتمندان»
    \item جنبش جلیقه‌زردها (۲۰۱۸)
    \item فاصله از طبقات محروم
    \item نخبه‌گرایی تکنوکراتیک
\end{itemize}

\textbf{سؤال:}
آیا ماکرون واقعاً «میانه» است، یا لیبرال راست‌میانه که خود را میانه می‌نامد؟ آیا «نه چپ، نه راست» موضعی جدید است، یا همان لیبرالیسم با نام جدید؟

\end{tcolorbox}

%═══════════════════════════════════════════════════════════════
\section{میانه در نظام‌های حزبی}
%═══════════════════════════════════════════════════════════════

\subsection{نظام دوحزبی}


در نظام‌های دوحزبی (مانند آمریکا و بریتانیا)، «میانه» جایی است که دو حزب بزرگ برای جذب رأی‌دهندگان بی‌طرف رقابت می‌کنند.

\textbf{نظریهٔ رأی‌دهندهٔ میانی:}
آنتونی داونز در \textit{نظریهٔ اقتصادی دموکراسی} (۱۹۵۷) استدلال کرد که در نظام دوحزبی، هر دو حزب به سمت میانه حرکت می‌کنند تا رأی‌دهندهٔ میانی را جذب کنند.

اما در دهه‌های اخیر، این نظریه به چالش کشیده شده:
\begin{itemize}
    \item قطبی‌شدن احزاب
    \item کاهش رأی‌دهندگان میانی
    \item اهمیت «پایگاه» حزبی
    \item رسانه‌های اجتماعی و حباب‌های اطلاعاتی
\end{itemize}

در آمریکای امروز، فاصلهٔ دموکرات‌ها و جمهوری‌خواهان بیشتر از هر زمان دیگری است.


\subsection{احزاب میانه در نظام‌های چندحزبی}

\begin{table}[H]
\centering
\begin{tabular}{|>{\columncolor{centergreen!10}}r|p{3cm}|p{4cm}|p{3cm}|}
\hline
\rowcolor{centergreen!30}
\textbf{کشور} & \textbf{حزب میانه} & \textbf{ایدئولوژی} & \textbf{وضعیت} \\
\hline
\textbf{آلمان} & FDP & لیبرال کلاسیک & ائتلافی \\
\hline
\textbf{بریتانیا} & لیبرال دموکرات‌ها & سوسیال-لیبرال & ضعیف \\
\hline
\textbf{فرانسه} & Renaissance (ماکرون) & لیبرال میانه & حاکم \\
\hline
\textbf{هلند} & D66 & لیبرال پیشرو & ائتلافی \\
\hline
\textbf{کانادا} & لیبرال‌ها & لیبرال میانه & حاکم \\
\hline
\end{tabular}
\caption{احزاب میانه در کشورهای مختلف}
\end{table}


احزاب میانه در نظام‌های چندحزبی اغلب نقش «پادشاه‌ساز» دارند: می‌توانند با چپ یا راست ائتلاف کنند و اهرم قدرت قابل توجهی دارند.

اما این احزاب با چالش‌هایی مواجه‌اند:
\begin{itemize}
    \item فشار از دو سو
    \item اتهام بی‌اصالتی
    \item دشواری تمایز هویتی
    \item قطبی‌شدن جامعه که میانه را می‌فشرد
\end{itemize}


%═══════════════════════════════════════════════════════════════
\section{نقد میانه‌روی}
%═══════════════════════════════════════════════════════════════

\begin{tcolorbox}[
    enhanced,
    colback=leftred!8,
    colframe=leftred,
    boxrule=1.5pt,
    arc=5pt,
    fonttitle=\bfseries,
    title={\rl{🔴 مناظره: میانه‌روی - فضیلت یا ضعف؟}}
]

\textbf{نقد از چپ:}
\begin{itemize}
    \item میانه‌روی = حفظ وضع موجود
    \item وضع موجود ناعادلانه است
    \item «میانه» بین ستم و عدالت، ستم است
    \item سانتریست‌ها با قدرت همراه‌اند
    \item تغییر واقعی رادیکال است
\end{itemize}



\textbf{نقد از راست:}
\begin{itemize}
    \item میانه‌روی = بی‌اصولی
    \item میانه‌روها اصول ثابت ندارند
    \item «میانه» در طول زمان به چپ می‌رود
    \item سانتریست‌ها همیشه تسلیم چپ می‌شوند
    \item ارزش‌ها قابل معامله نیستند
\end{itemize}


\vspace{3mm}
\textbf{دفاع از میانه‌روی:}
\begin{itemize}
    \item افراط همیشه خطرناک است، چپ یا راست
    \item تغییر تدریجی پایدارتر از انقلاب است
    \item واقع‌گرایی فضیلت است، نه ضعف
    \item دموکراسی نیازمند مصالحه است
\end{itemize}
\end{tcolorbox}

\begin{tcolorbox}[
    enhanced,
    colback=centergreen!10,
    colframe=centergreen,
    boxrule=1.5pt,
    arc=5pt,
    fonttitle=\bfseries,
    title={\rl{🟢 نقل‌قول: مارتین لوتر کینگ دربارهٔ میانه‌روها}}
]
\textit{«باید اعتراف کنم که در این چند سال گذشته، سخت از سفیدپوست میانه‌رو ناامید شده‌ام. تقریباً به این نتیجه رسیده‌ام که سنگ بزرگ بر سر راه نگرو به سوی آزادی، نه عضو کو کلاکس کلان یا شورای شهروندان سفید، بلکه سفیدپوست میانه‌رویی است که بیشتر به 'نظم' وفادار است تا به عدالت.»}

\hfill --- مارتین لوتر کینگ، \textit{نامه از زندان بیرمنگام}، ۱۹۶۳
\end{tcolorbox}

%═══════════════════════════════════════════════════════════════
\section{پنجرهٔ اورتون: میانهٔ متحرک}
%═══════════════════════════════════════════════════════════════

\begin{tcolorbox}[
    enhanced,
    colback=rightblue!8,
    colframe=rightblue,
    boxrule=1.5pt,
    arc=5pt,
    fonttitle=\bfseries,
    title={\rl{🔵 مفهوم کلیدی: پنجرهٔ اورتون}}
]
\textbf{پنجرهٔ اورتون} (Overton Window) مفهومی است که توسط جوزف اورتون (۱۹۶۰-۲۰۰۳) مطرح شد. این پنجره نشان‌دهندهٔ محدودهٔ ایده‌هایی است که در یک زمان معین «قابل قبول» یا «متعارف» تلقی می‌شوند.

ایده‌ها می‌توانند در طیفی قرار گیرند:
\begin{center}
غیرقابل تصور ← رادیکال ← قابل قبول ← معقول ← رایج ← سیاست ← رایج ← معقول ← قابل قبول ← رادیکال ← غیرقابل تصور
\end{center}

آنچه در «پنجره» قرار دارد، میانهٔ سیاسی آن زمان است. اما این پنجره می‌تواند جابجا شود.
\end{tcolorbox}

\begin{center}
\begin{tikzpicture}[scale=0.9, every node/.style={font=\small}]

% Main axis
\draw[line width=2pt, color=darkgray] (0,0) -- (14,0);
\node[below] at (0,0) {\rl{چپ افراطی}};
\node[below] at (14,0) {\rl{راست افراطی}};

% 1950s window
\draw[line width=3pt, color=rightblue] (4,1) -- (9,1);
\draw[line width=1pt, color=rightblue, dashed] (4,0) -- (4,1);
\draw[line width=1pt, color=rightblue, dashed] (9,0) -- (9,1);
\node[above, color=rightblue] at (6.5,1.2) {\rl{پنجره ۱۹۵۰}};

% 2020s window (shifted left)
\draw[line width=3pt, color=centergreen] (2,2.5) -- (7,2.5);
\draw[line width=1pt, color=centergreen, dashed] (2,0) -- (2,2.5);
\draw[line width=1pt, color=centergreen, dashed] (7,0) -- (7,2.5);
\node[above, color=centergreen] at (4.5,2.7) {\rl{پنجره ۲۰۲۰}};

% Arrow showing shift
\draw[->, line width=2pt, color=goldaccent] (6.5,1.5) -- (4.5,2.2);
\node[right, color=goldaccent] at (5.5,1.85) {\rl{جابجایی}};

% Title
\node[above, font=\bfseries\large] at (7,4) {\rl{جابجایی پنجرهٔ اورتون در طول زمان}};

\end{tikzpicture}
\end{center}


\textbf{مثال: ازدواج همجنس‌گرایان}

\begin{itemize}
    \item \textbf{۱۹۹۰:} غیرقابل تصور
    \item \textbf{۲۰۰۰:} رادیکال
    \item \textbf{۲۰۱۰:} قابل قبول
    \item \textbf{۲۰۱۵:} قانونی در آمریکا
    \item \textbf{۲۰۲۰:} موضع میانه
\end{itemize}

کسی که در ۱۹۹۰ از ازدواج همجنس‌گرایان دفاع می‌کرد، «چپ افراطی» بود. امروز همان موضع «میانه» است.

\textbf{نتیجه:} میانه ثابت نیست. آنچه امروز میانه است، ممکن است فردا راست یا چپ باشد. این نسبی بودن میانه، پرسش‌های اخلاقی مهمی مطرح می‌کند: آیا باید همیشه در میانه بود، حتی اگر میانه ناعادلانه باشد؟


%═══════════════════════════════════════════════════════════════
\section{آیندهٔ میانه}
%═══════════════════════════════════════════════════════════════

\subsection{بحران احزاب میانه}


در دهه‌های اخیر، احزاب میانه (سوسیال دموکرات و محافظه‌کار میانه) در بسیاری از کشورها تضعیف شده‌اند:

\textbf{شواهد:}
\begin{itemize}
    \item کاهش آرای احزاب سنتی
    \item رشد احزاب پوپولیست چپ و راست
    \item کاهش عضویت حزبی
    \item بی‌اعتمادی به نخبگان سیاسی
\end{itemize}

\textbf{دلایل احتمالی:}
\begin{itemize}
    \item نابرابری فزاینده
    \item احساس «رها شدن» طبقات متوسط و پایین
    \item شکست سیاست‌های میانه در حل مشکلات
    \item رسانه‌های اجتماعی و قطبی‌شدن
    \item بحران‌های مالی ۲۰۰۸ و کرونا
\end{itemize}


\subsection{آیا میانه بازمی‌گردد؟}


\textbf{سناریوی بازگشت:}
\begin{itemize}
    \item خستگی از قطبی‌شدن
    \item شکست پوپولیسم در حکومت
    \item نیاز به مصالحه برای حل مشکلات (اقلیم، همه‌گیری)
    \item نسل جدید سیاستمداران میانه
\end{itemize}

\textbf{سناریوی ادامهٔ قطبی‌شدن:}
\begin{itemize}
    \item مشکلات ساختاری حل نشده
    \item الگوریتم‌های رسانه‌های اجتماعی
    \item هویت‌گرایی فزاینده
    \item بی‌اعتمادی عمیق‌تر
\end{itemize}

شاید آینده نه بازگشت به میانهٔ قدیم، بلکه بازتعریف میانه باشد. میانهٔ جدید ممکن است ترکیبی از عناصر مختلف باشد که در چارچوب سنتی چپ-راست نمی‌گنجد.


%═══════════════════════════════════════════════════════════════
\section*{خلاصهٔ فصل}
%═══════════════════════════════════════════════════════════════

\begin{tcolorbox}[
    enhanced,
    colback=lightbg,
    colframe=darkgray,
    boxrule=2pt,
    arc=8pt,
    fonttitle=\bfseries\large,
    title={\rl{📝 خلاصهٔ فصل چهاردهم}}
]
\begin{itemize}
    \item «میانه» سه معنا دارد: \textbf{نقطه} (موضع محتوایی)، \textbf{روش} (پراگماتیسم)، و \textbf{فضای اجتماعی} (جای اکثریت).
    
    \item میانه‌روی ریشه در سنت فلسفی دارد: \textbf{حد وسط ارسطو}، \textbf{اعتدال دینی}، و \textbf{پراگماتیسم آمریکایی}.
    
    \item «سانتریسم رادیکال» تلاش می‌کند میانه‌روی را با اصلاحات بنیادین ترکیب کند. \textbf{ماکرون} نمونه‌ای معاصر است.
    
    \item منتقدان از چپ میانه‌روی را \textbf{حفظ وضع موجود} و منتقدان از راست آن را \textbf{بی‌اصولی} می‌دانند.
    
    \item \textbf{پنجرهٔ اورتون} نشان می‌دهد که میانه ثابت نیست؛ با تغییر افکار عمومی جابجا می‌شود.
    
    \item احزاب میانه در \textbf{بحران} هستند: قطبی‌شدن جوامع، فضای میانه را می‌فشرد.
    
    \item آیندهٔ میانه نامعلوم است: یا بازگشت، یا ادامهٔ قطبی‌شدن، یا بازتعریف.
\end{itemize}
\end{tcolorbox}

%═══════════════════════════════════════════════════════════════
\section*{پرسش‌های تأملی}
%═══════════════════════════════════════════════════════════════

\begin{tcolorbox}[
    enhanced,
    colback=centergreen!8,
    colframe=centergreen,
    boxrule=1.5pt,
    arc=5pt,
    fonttitle=\bfseries,
    title={\rl{❓ پرسش‌هایی برای تأمل و بحث}}
]
\begin{enumerate}
    \item آیا «حق» همیشه در میانه است؟ مثال‌هایی بزنید که میانه‌روی درست بود، و مثال‌هایی که موضع رادیکال درست بود.
    
    \item نقد کینگ از میانه‌روهای سفید چه درس‌هایی برای امروز دارد؟ آیا گاهی میانه‌روی همدستی با ظلم است؟
    
    \item اگر پنجرهٔ اورتون در طول زمان تغییر می‌کند، چرا باید در میانه بود؟ آیا نباید تلاش کرد پنجره را جابجا کرد؟
    
    \item آیا ماکرون واقعاً «میانه» است، یا لیبرال راست‌میانه؟ تفاوت کجاست؟
    
    \item چرا احزاب میانه در بسیاری از کشورها تضعیف شده‌اند؟ آیا این روند برگشت‌پذیر است؟
    
    \item آیا پراگماتیسم («هر چه کار می‌کند») یک فلسفهٔ سیاسی کافی است، یا نیازمند اصول اخلاقی مستقل هستیم؟
\end{enumerate}
\end{tcolorbox}

\vspace{10mm}
\begin{center}
\rule{0.6\textwidth}{1pt}

\vspace{3mm}
\textit{«اگر در میانهٔ جاده بایستی، از هر دو طرف زیر می‌گیرند.»}

\vspace{2mm}
--- آنورین بِوان، سیاستمدار بریتانیایی

\vspace{3mm}
\rule{0.6\textwidth}{1pt}
\end{center}